\subsection*{Overview}
Brain signals are fundamental in neuroengineering. This course covers methods for analyzing and interpreting brain signals across scales—from single neurons to whole-brain activity. Students will learn how these signals are generated, the main acquisition methodologies, and how to analyze, interpret, and discuss neural data.
The course provides a critical understanding of neuronal signal analysis, from single-spike characterization to multi-unit activity and EEG signal processing. Students will acquire the skills to manipulate, analyze, and interpret electrophysiological data.\\ 


\textbf{Aim 1:} Understanding the neurophysiological bases and generation mechanisms of brain signals.

- Describe how different neural signals are generated across spatio-temporal scales (micro, meso, macro).  

- Distinguish between brain signal types and their recording methodologies.

- \textit{Assessed during the oral examination through critical discussion of the methods presented.} 

\textbf{Aim 2:} Extracting information from brain signals.
 
- Identify optimal pre-processing and feature extraction strategies.  

- Apply and critically evaluate algorithms for signal analysis.

- \textit{Assessed through the group-based assignments.} 


\textbf{Aim 3:} Solving real-case problems in neural signal analysis.

- Apply course techniques to real-case data problems in working groups.  

- Collaborate in team-based lab simulations with role differentiation.

- \textit{Assessed through the group-based assignments.} \\ 


Lectures, flipped classrooms, problem-based learning, group work, blended learning.\\

\paragraph*{Prerequisites}
Basic Matlab programming, neurophysiology, physics, linear algebra.\\

\paragraph*{Recommended Reading/Bibliography}
\href{https://direct.mit.edu/books/monograph/4013/Analyzing-Neural-Time-Series-DataTheory-and}{Analyzing Neural Time Series Data: Theory and Practice}

\subsection*{Exam Description}
The exam consists of project assignments (group work) and an oral exam. Students will self-organize into small groups (max 3) and complete several activities during the semester. Each assignment will be evaluated for completeness and overall quality.  
Assignments start in mid-May, are developed in Python, and presented at the end of the course.

\paragraph*{Grading}
Final grade = oral exam grade + up to 2 additional points from the assignment.